\section{Introduction}
Economic expectations matter before they become survey responses or numerical forecasts. A newspaper may report a business leader anticipating recovery, an economist expecting unemployment to rise, or an editorial predicting a decline in prices. Such statements can preserve the direction, uncertainty, and attribution of expectations at a particular date. Their historical value is considerable, but turning them into comparable observations requires more than locating future-tense language. A statement about slower growth differs from a prediction of falling output; a forecast about jobs differs from a forecast about the unemployment rate; and a broad outlook may not specify an endpoint that can be scored mechanically.

Language models make it feasible to generate structured candidates from large text collections at relatively low marginal API cost. Whether those candidates support economic conclusions is a separate empirical question. Cheap extraction can be useful even when imperfect, but its usefulness depends on the errors it introduces, the evidence available for validating them, and the outcome rule applied downstream. We therefore examine three linked stages: extraction, source interpretation, and evaluation against economic outcomes.

The main application uses the annual ProQuest economy exports for 1900--2010. We retain 274,070 machine-verified candidates without automatically deduplicating identical stripped labels. The primary historical analysis uses only these annual exports; the raw annual candidates form a sensitivity set. A separate Library of Congress (LOC) experiment supplies development evidence about extraction cost and reference agreement. Its pages and economic results never enter the ProQuest estimation sample. A separately authorized 2011--2026 ProQuest extension contains 91,118 retained candidates and evaluates 56,886 mature, supported claims. We report its exploratory transfer results in the main text, with full details in Appendix~\ref{app:modern}; historical and modern estimates remain separate.

The paper makes three contributions. First, it documents a concrete computational cost--quality comparison, including the precision--recall consequences of verification. The most economical illustrated two-stage configuration costs approximately one-fifteenth as much as the higher-recall configuration in the recorded page-scaling exercise. This establishes a development tradeoff, not a universal price advantage or an end-to-end cost estimate for ProQuest. Second, it reports source-review evidence with the sampling frames, field denominators, and reviewer provenance made explicit. Third, it demonstrates that reasonable-looking changes to benchmark definitions can materially alter conclusions about the same historical claims.

The last contribution is the principal economic finding. Under the original employment rule, the matched sample is correct 34.2\% of the time and exceeds an always-improve baseline by 14.2 percentage points. Under sign-only scoring, accuracy reaches 57.6\%, but the advantage is 10.8 points and its block interval includes zero. Replacing falling unemployment with rising payroll employment changes the baseline-relative point estimate from positive to negative. Higher raw accuracy therefore does not, by itself, identify a better interpretation of what the newspaper predicted.

This argument concerns rules implemented in this repository. We do not assert that other researchers routinely use an undocumented defective threshold, nor infer a field-wide failure from a hypothetical example. The manuscript also makes no new mathematical priority claim: the definitions and elementary identities below organize the empirical comparison. The evidence is an observed dependence of economic conclusions on specific, auditable measurement choices.

Predictive models provide an application rather than the criterion for the paper's success. A recurrent model can distinguish correct from incorrect machine-scored claims without proving that the extraction is valid, that recurrence supplies the improvement, or that the model forecasts recessions. A historical aggregate comparison shows a small average gain from newspaper features, but that gain is uncertain and reverses when 2007--2009 evaluation months are omitted. Reporting these limitations is necessary to establish what this dataset currently measures.

